\documentclass[11pt]{article}
\usepackage[margin=1in]{geometry}
\usepackage[T1]{fontenc}
\usepackage{lmodern}
\usepackage{amsmath,amssymb,amsthm,mathtools}
\usepackage{graphicx}
\usepackage[hidelinks]{hyperref}
\usepackage{microtype}
\usepackage{enumitem}

\newtheorem{theorem}{Theorem}
\newtheorem{corollary}[theorem]{Corollary}
\newtheorem{lemma}[theorem]{Lemma}
\newtheorem{proposition}[theorem]{Proposition}

\newcommand{\F}{\mathcal F}
\newcommand{\Lip}{\operatorname{Lip}}
\newcommand{\Lipz}{\operatorname{Lip}_0}
\newcommand{\R}{\mathbb R}
\newcommand{\D}{\mathcal D}
\newcommand{\dist}{\operatorname{dist}}
\newcommand{\diver}{\operatorname{div}}
\newcommand{\clspan}{\overline{\operatorname{span}}}

\begin{document}

\begin{center}
{\Large Full solution packet for arXiv:2108.02727}\\[4pt]
{\large Collapsing subsets in Lipschitz-free spaces and an explicit flux model for persistence diagrams}
\end{center}

\paragraph{Status.}
\textbf{full\_solution\_likely\_valid}.  This packet answers the concrete
problem in the conclusion of Chad Giusti and Darrick Lee,
\emph{Signatures, Lipschitz-free spaces, and paths of persistence diagrams},
arXiv:2108.02727.  It first gives the exact functorial answer for every metric
pair, then combines it with the Euclidean representation of
C\'uth--Kalenda--Kaplick\'y to obtain an explicit vector-field representation
for ordinary persistence diagrams.

\paragraph{Source question (Conclusion, p.~26).}
The source asks for ``an explicit description of the Lipschitz-free space of
the quotient metric space of a metric pair'' and, specifically, whether known
representations of $\F(\R^N)$ can be adapted to represent persistence
diagrams.

\begin{center}
\includegraphics[width=.96\linewidth]{figures/open_problem_crop.png}
\end{center}

\section*{Result at a glance}

Let $X/A$ denote the pointed metric space obtained by collapsing a nonempty
closed subset $A$ of a metric space $X$ to one point.  With $X$ based at a
point of $A$, put
\[
  \F_X(A)=\clspan\{\delta_X(a):a\in A\}\subseteq \F(X).
\]
There is a canonical linear isometry
\[
  \boxed{\quad \F(X/A)\cong \F(X)/\F_X(A).\quad}                 \tag{1}
\]
Thus every concrete model of $\F(X)$ automatically induces one for the
collapsed metric pair: quotient the model by the vectors representing $A$.

For ordinary persistence diagrams, rotate the birth--death half-plane so that
\[
 H=\R\times(0,\infty),\qquad \partial H=\R\times\{0\}.
\]
The diagonal becomes $\partial H$.  Then (1) and the theorem of
C\'uth--Kalenda--Kaplick\'y give the explicit isometric representation
\[
 \boxed{
 \F(\overline H/\partial H)
 \cong
 \frac{L^1(H;\R^2)}
 {\{g\in L^1(H;\R^2):\diver g=0\text{ in }\D'(H)\}} .}         \tag{2}
\]
The fact that divergence is taken \emph{inside $H$}, rather than after
zero-extension to $\R^2$, is the key boundary effect: flux may disappear
through the collapsed diagonal.

If $D,E$ are finite persistence diagrams and $\nu_D,\nu_E$ are their counting
measures after rotation into $H$, then (2) gives the exact flux formula
\[
 \boxed{
 W_1^\partial(D,E)
 =\inf\left\{\|g\|_{L^1(H)}:
       g\in L^1(H;\R^2),\ \diver g=\nu_E-\nu_D\text{ in }\D'(H)
       \right\}.}                                             \tag{3}
\]

\section*{Proof intuition}

There are two simple dual facts behind the result.

First, a Lipschitz function on $X/A$ is exactly a Lipschitz function on $X$
that vanishes on $A$.  The two Lipschitz constants agree: besides the direct
bound through $d(x,y)$, vanishing on $A$ gives the second bound through
$d(x,A)+d(y,A)$.  Since Lipschitz-free spaces are the canonical preduals of
these function spaces, collapsing $A$ amounts to quotienting out the closed
span of the evaluation vectors on $A$.

Second, in the Euclidean representation a vector field $g$ acts on a
Lipschitz function $f$ by $\int g\cdot\nabla f$.  After the boundary is
collapsed, the only test functions are those with zero boundary values.
Consequently, the invisible vector fields are precisely those whose
divergence vanishes in the interior.  Boundary flux is deliberately not
tested.  An atom of a persistence diagram is a unit source whose flux can
either travel to another atom or exit through the boundary; minimizing the
$L^1$ cost gives exactly partial $W_1$.

\section*{The general metric-pair theorem}

For a nonempty closed $A\subseteq X$, write
\[
 d_A([x],[y])=\min\{d(x,y),d(x,A)+d(y,A)\}
\]
for the quotient metric, with the evident interpretation when one class is
the collapsed point $*$.  Choose $a_0\in A$ as the base point of $X$.

\begin{theorem}[Free spaces commute with collapsing]\label{thm:collapse}
Let $X$ be a metric space and let $A\subseteq X$ be nonempty and closed.
The quotient map $q:X\to X/A$ induces a canonical surjective linear map
\[
 Q:\F(X)\longrightarrow \F(X/A),\qquad Q\delta_X(x)=\delta_{X/A}(q(x)),
\]
and
\[
 \ker Q=\F_X(A).
\]
The induced map $\F(X)/\F_X(A)\to\F(X/A)$ is a linear isometry.
\end{theorem}

\begin{proof}
Let
\[
 I_A=\{f\in\Lipz(X): f|_A=0\}.
\]
Composition with $q$ maps $\Lipz(X/A)$ into $I_A$.  This map is onto: if
$f\in I_A$, define $\widetilde f([x])=f(x)$ and $\widetilde f(*)=0$.
It is also isometric.  Indeed, if $L=\|f\|_{\Lip}$, then
\[
 |f(x)-f(y)|\le Ld(x,y)
\]
and, because $f$ vanishes on $A$,
\[
 |f(x)-f(y)|\le |f(x)|+|f(y)|
 \le L\bigl(d(x,A)+d(y,A)\bigr).
\]
Taking the smaller bound shows that $\widetilde f$ is $L$-Lipschitz on
$X/A$.  The reverse norm inequality follows because $q$ is $1$-Lipschitz.
Therefore
\[
 \Lipz(X/A)\xrightarrow[cong]{\ f\mapsto f\circ q\ } I_A.    \tag{4}
\]

The annihilator of $\F_X(A)$ in $\F(X)^*=\Lipz(X)$ is exactly $I_A$.
The adjoint of the induced map
\[
 \overline Q:\F(X)/\F_X(A)\longrightarrow\F(X/A)
\]
is precisely (4), hence is an onto isometry.  It follows directly that
$\overline Q$ is an isometry: for $u\in\F(X)/\F_X(A)$,
\[
 \|\overline Qu\|
 =\sup_{\|h\|_{\Lip}\le1}|\langle u,h\circ q\rangle|=\|u\|.
\]
Its range is closed, and its annihilator is $\ker \overline Q^*=\{0\}$, so
the range is also dense.  Thus $\overline Q$ is onto.  The asserted kernel
identity follows.
\end{proof}

\paragraph{Interpretation.}
The theorem is base-point independent in the usual canonical sense.  The
choice $a_0\in A$ merely makes every $\delta_X(a)$ unambiguous.  Also,
$\F_X(A)$ is canonically the embedded copy of $\F(A)$; the isometry of that
embedding follows, for example, by the scalar McShane extension theorem.

\section*{Euclidean metric pairs}

We record the exact mechanism for adapting the known Euclidean model.  Let
$U\subset\R^n$ be nonempty, open, and convex, and let $A$ be a nonempty
closed subset of $\overline U$.  Define
\[
 \mathcal Y_A=
 \left\{g\in L^1(U;\R^n):
   \int_U g\cdot\nabla f=0
   \quad\text{for every }f\in\Lip(\overline U)\text{ with }f|_A=0
 \right\}.                                                     \tag{5}
\]

\begin{proposition}[Explicit model for a Euclidean pair]\label{prop:pair}
There is a canonical linear isometry
\[
 \F(\overline U/A)\cong L^1(U;\R^n)/\mathcal Y_A.              \tag{6}
\]
Under (6), $g$ acts on $f\in\Lip(\overline U)$ with $f|_A=0$ by
$\int_U g\cdot\nabla f$.
\end{proposition}

\begin{proof}
C\'uth--Kalenda--Kaplick\'y, Theorem~1.1, identify $\F(U)$ isometrically
with a quotient of $L^1(U;\R^n)$ via the pairing
\[
 (g,f)\longmapsto\int_U g\cdot\nabla f.                        \tag{7}
\]
Passing from $U$ to its completion $\overline U$ does not change the free
space: the evaluation vectors from the dense subset $U$ have dense span and
the metric embedding extends continuously.  Changing the base point only
adds constants to dual Lipschitz functions, so it does not alter the gradient
pairing (7).  Apply
Theorem~\ref{thm:collapse}.  By (4), the dual functions that remain after
collapsing $A$ are exactly the Lipschitz functions vanishing on $A$.
Therefore the preimage, under (7), of the subspace being quotiented out is
exactly the annihilator (5).  Taking the quotient gives (6), with no change
of norm.
\end{proof}

For comparison, before collapsing a boundary the kernel in the
C\'uth--Kalenda--Kaplick\'y theorem consists of fields whose zero-extension
has divergence zero on all of $\R^n$.  Formula (5) shows exactly how a
collapsed subset changes that condition.

\section*{The persistence half-plane}

Let $H=\R\times(0,\infty)$ and $A=\partial H$.  Distributional divergence
in $H$ is defined by
\[
 \langle\diver g,\varphi\rangle=-\int_H g\cdot\nabla\varphi
 \qquad(\varphi\in C_c^\infty(H)).
\]

\begin{lemma}[The collapsed-boundary kernel]\label{lem:boundary}
For the pair $(\overline H,\partial H)$, the subspace in (5) is
\[
 \mathcal Y_{\partial H}
 =\mathcal Z_H
 :=\{g\in L^1(H;\R^2):\diver g=0\text{ in }\D'(H)\}.           \tag{8}
\]
\end{lemma}

\begin{proof}
If $g\in\mathcal Y_{\partial H}$, test (5) with
$\varphi\in C_c^\infty(H)$ to obtain $\diver g=0$ in $H$.

Conversely, suppose $\diver g=0$ in $H$, and let
$f\in\Lip(\overline H)$ vanish on $\partial H$.  Extend $f$ by zero to the
lower half-plane.  The extension $F$ is Lipschitz on $\R^2$: for
$(s,t)\in H$, $|f(s,t)|\le \|f\|_{\Lip}t$, which also handles pairs lying on
opposite sides of the boundary.

Choose $\chi_R\in C_c^\infty(\R^2)$ with $\chi_R=1$ on $B(0,R)$,
$\chi_R=0$ outside $B(0,2R)$, and $|\nabla\chi_R|\le C/R$.  Put
$F_R=\chi_RF$.  Translate $F_R$ upward by $\varepsilon$ and mollify at a
scale smaller than $\varepsilon/2$.  This produces functions in
$C_c^\infty(H)$ whose gradients converge almost everywhere to $\nabla F_R$
and are bounded uniformly in $L^\infty$.  Pairing with $g\in L^1$ and using
dominated convergence gives
\[
 \int_H g\cdot\nabla F_R=0.                                   \tag{9}
\]
As $R\to\infty$, $\nabla F_R\to\nabla f$ almost everywhere on $H$.
Moreover the gradients are uniformly bounded: on the annulus where
$\nabla\chi_R\ne0$,
\[
 |F|\,|\nabla\chi_R|
 \le \|f\|_{\Lip}t\,C/R\le 2C\|f\|_{\Lip}.
\]
Dominated convergence in (9) now yields $\int_Hg\cdot\nabla f=0$.
Thus $g\in\mathcal Y_{\partial H}$.
\end{proof}

\begin{corollary}[Vector-field model for persistence diagrams]\label{cor:pd}
Let
\[
 \Omega=\{(b,d)\in\R^2:b\le d\},\qquad
 \Delta=\{(t,t):t\in\R\},
\]
with the Euclidean metric.  Then
\[
 \F(\Omega/\Delta)
 \cong L^1(H;\R^2)/\mathcal Z_H                              \tag{10}
\]
canonically and isometrically, after the rotation
\[
 \rho(b,d)=\left(\frac{b+d}{\sqrt2},\frac{d-b}{\sqrt2}\right).
\]
For finite diagrams $D,E$, let
$\nu_D=\sum_{z\in D}\delta_{\rho(z)}$ and define $\nu_E$ similarly, omitting
diagonal points.  The class of $D-E$ consists exactly of the fields $g$ with
\[
 \diver g=\nu_E-\nu_D\quad\text{in }\D'(H),                  \tag{11}
\]
and its norm is the infimum of $\|g\|_{L^1(H)}$ over (11).  Equivalently,
formula (3) holds.
\end{corollary}

\begin{proof}
The rotation $\rho$ is an isometry from $\Omega$ onto $\overline H$ and maps
$\Delta$ onto $\partial H$.  Proposition~\ref{prop:pair} and
Lemma~\ref{lem:boundary} give (10).

The pairing convention gives
\[
 \int_H g\cdot\nabla f
 =-\langle\diver g,f\rangle.
\]
Hence (11) represents
$\langle\nu_D-\nu_E,f\rangle$, exactly the evaluation vector of $D-E$.
Two representatives differ precisely by an element of $\mathcal Z_H$, and
the quotient norm is the stated infimum.  Finally, Giusti--Lee define
$\F(\Omega/\Delta)$ as the completion of finite virtual diagrams in the
partial $W_1$ norm.  Therefore that quotient norm equals
$W_1^\partial(D,E)$.
\end{proof}

For a single pair $x,y\in H$, Corollary~\ref{cor:pd} specializes to the
transparent identity
\[
 \inf_{\diver g=\delta_y-\delta_x}\|g\|_1
 =\min\{|x-y|,\dist(x,\partial H)+\dist(y,\partial H)\}.
\]
The two terms are the two possible flux routes: directly from $x$ to $y$, or
from $x$ to the collapsed boundary and back out to $y$.

\section*{Verification and scope}

\paragraph{Checks performed.}
The quotient metric used in Theorem~\ref{thm:collapse} was checked against
the source definition.  The norm equality in (4) uses both branches of that
minimum, so it is not merely a contractive identification.  The supporting
paper's Theorem~1.1 was checked directly, including its warning that its
uncollapsed kernel uses zero-extension and divergence in all of $\R^n$.
Lemma~\ref{lem:boundary} independently derives the changed boundary
condition.  The signs in (11) follow from
$\langle\diver g,f\rangle=-\int g\cdot\nabla f$.

\paragraph{What ``explicit'' means here.}
Formula (10) is an exact, isometric Banach-space quotient, and (11) gives a
concrete PDE description of every diagram vector.  It is not a unique
coordinate representation: adding a divergence-free field does not change
the element.  Nor is it finite-dimensional.  Those limitations already
occur in the known representation of $\F(\R^2)$ that motivated the source
question.

\paragraph{Novelty assessment.}
Exact-phrase and structural searches were run for the quotient identity, a
collapsed-subset representation, and a divergence/flux representation of
persistence diagrams.  The searches recovered the ambient convex-domain
theorem used above and later extensions of ambient free-space
representations, but no explicit collapsed-diagonal formula (10) or diagram
flux formula (3).  The abstract identity (1) is elementary enough that it may
be folklore; the claimed mathematical contribution of this packet is the
complete derivation and its persistence-diagram specialization.  Novelty
confidence is therefore moderate, while proof confidence is high.

\paragraph{Human-review recommendation.}
Review as a full answer to the concrete problem on p.~26.  The two points
most worth checking are the weak-$*$ approximation in
Lemma~\ref{lem:boundary} and the interpretation of partial $W_1$ in
Corollary~\ref{cor:pd}.  The result concerns the $p=1$ geometry used by
Lipschitz-free spaces; it does not claim analogous linearizations for
$p>1$.

\begin{thebibliography}{9}
\bibitem{GiustiLee}
C. Giusti and D. Lee,
\emph{Signatures, Lipschitz-free spaces, and paths of persistence diagrams},
arXiv:2108.02727, especially the Conclusion, p.~26.

\bibitem{CuthKalendaKaplicky}
M. C\'uth, O. F. K. Kalenda, and P. Kaplick\'y,
\emph{Isometric representation of Lipschitz-free spaces over convex domains
in finite-dimensional spaces},
Mathematika 63 (2017), 538--552; arXiv:1610.03966,
Theorem~1.1.
\end{thebibliography}

\end{document}
